\documentclass[10pt,a4paper]{amsart}
\usepackage{amsmath,amssymb,amsthm,graphicx}
% Retained preamble for the Stacks Project chapter morphisms.tex (Morphisms of Schemes).
% The chapter file pinned-original-source.tex carries only \input{preamble} on its first line, so no
% preamble text lives inside the chapter itself.  The Stacks Project uses one single shared preamble for
% every chapter; the snapshot stages it verbatim in sources/stacks-chow/preamble.tex.  The block below is
% copied mechanically, line for line, from that file's lines 14-18 (the Xy-pic setup) and lines 154-217
% (the native theorem declarations together with the mathematical macros).  Nothing was added, removed or
% redefined.  The wrapper class and the amsmath/amssymb/amsthm/graphicx base packages are supplied by the
% assembler, so the class/\externaldocument part of the shared preamble is not repeated here.
% For commutative diagrams we use Xy-pic
\usepackage[all]{xy}

% We use 2cell for 2-commutative diagrams.
\xyoption{2cell}
% Theorem environments.
%
\theoremstyle{plain}
\newtheorem{theorem}[subsection]{Theorem}
\newtheorem{proposition}[subsection]{Proposition}
\newtheorem{lemma}[subsection]{Lemma}

\theoremstyle{definition}
\newtheorem{definition}[subsection]{Definition}
\newtheorem{example}[subsection]{Example}
\newtheorem{exercise}[subsection]{Exercise}
\newtheorem{situation}[subsection]{Situation}

\theoremstyle{remark}
\newtheorem{remark}[subsection]{Remark}
\newtheorem{remarks}[subsection]{Remarks}

\numberwithin{equation}{subsection}

% Macros
%
\def\lim{\mathop{\mathrm{lim}}\nolimits}
\def\colim{\mathop{\mathrm{colim}}\nolimits}
\def\Spec{\mathop{\mathrm{Spec}}}
\def\Hom{\mathop{\mathrm{Hom}}\nolimits}
\def\Ext{\mathop{\mathrm{Ext}}\nolimits}
\def\SheafHom{\mathop{\mathcal{H}\!\mathit{om}}\nolimits}
\def\SheafExt{\mathop{\mathcal{E}\!\mathit{xt}}\nolimits}
\def\Sch{\mathit{Sch}}
\def\Mor{\mathop{\mathrm{Mor}}\nolimits}
\def\Ob{\mathop{\mathrm{Ob}}\nolimits}
\def\Sh{\mathop{\mathit{Sh}}\nolimits}
\def\NL{\mathop{N\!L}\nolimits}
\def\CH{\mathop{\mathrm{CH}}\nolimits}
\def\proetale{{pro\text{-}\acute{e}tale}}
\def\etale{{\acute{e}tale}}
\def\QCoh{\mathit{QCoh}}
\def\Ker{\mathop{\mathrm{Ker}}}
\def\Im{\mathop{\mathrm{Im}}}
\def\Coker{\mathop{\mathrm{Coker}}}
\def\Coim{\mathop{\mathrm{Coim}}}

% Boxtimes
%
\DeclareMathSymbol{\boxtimes}{\mathbin}{AMSa}{"02}

%
% Macros for moduli stacks/spaces
%
\def\QCohstack{\mathcal{QC}\!\mathit{oh}}
\def\Cohstack{\mathcal{C}\!\mathit{oh}}
\def\Spacesstack{\mathcal{S}\!\mathit{paces}}
\def\Quotfunctor{\mathrm{Quot}}
\def\Hilbfunctor{\mathrm{Hilb}}
\def\Curvesstack{\mathcal{C}\!\mathit{urves}}
\def\Polarizedstack{\mathcal{P}\!\mathit{olarized}}
\def\Complexesstack{\mathcal{C}\!\mathit{omplexes}}
% \Pic is the operator that assigns to X its picard group, usage \Pic(X)
% \Picardstack_{X/B} denotes the Picard stack of X over B
% \Picardfunctor_{X/B} denotes the Picard functor of X over B
\def\Pic{\mathop{\mathrm{Pic}}\nolimits}
\def\Picardstack{\mathcal{P}\!\mathit{ic}}
\def\Picardfunctor{\mathrm{Pic}}
\def\Deformationcategory{\mathcal{D}\!\mathit{ef}}

% --- appended 2026-10-12 (Chief): the project own definitions of the reference/slogan/history
% display environments, copied mechanically from the shared Stacks preamble lines 8-12
% (\usepackage{verbatim} plus the three \newenvironment lines, which the chapter preamble keeps
% verbatim from the project preamble). Nothing added, removed or redefined: these environments
% carry no mathematical content and the project itself typesets them as comments. Without these
% three lines a counted statement whose own text uses them fails with Environment slogan undefined.
\usepackage{verbatim}
\newenvironment{reference}{\comment}{\endcomment}
\newenvironment{slogan}{\comment}{\endcomment}
\newenvironment{history}{\comment}{\endcomment}
\usepackage{cleveref}
\crefname{originaldefinition}{original definition}{original definitions}
\Crefname{originaldefinition}{Original definition}{Original definitions}
\makeatletter
\newlabel{constructions-definition-projective-space}{{01NF}{source}{}{sourcepointer-1}{}}
\newlabel{constructions-definition-projective-space@cref}{{[originaldefinition][0][]01NF}{[1][0][]source}{}{}{}}
\makeatother
\crefname{originallemma}{original lemma}{original lemmas}
\Crefname{originallemma}{Original lemma}{Original lemmas}
\makeatletter
\newlabel{constructions-lemma-projective-space}{{01NE}{source}{}{sourcepointer-2}{}}
\newlabel{constructions-lemma-projective-space@cref}{{[originallemma][0][]01NE}{[1][0][]source}{}{}{}}
\makeatother
\crefname{originallemma}{original lemma}{original lemmas}
\Crefname{originallemma}{Original lemma}{Original lemmas}
\makeatletter
\newlabel{constructions-lemma-segre-embedding}{{01WD}{source}{}{sourcepointer-3}{}}
\newlabel{constructions-lemma-segre-embedding@cref}{{[originallemma][0][]01WD}{[1][0][]source}{}{}{}}
\makeatother
\crefname{originallemma}{original lemma}{original lemmas}
\Crefname{originallemma}{Original lemma}{Original lemmas}
\makeatletter
\newlabel{lemma-immersion-permanence}{{07RK}{source}{}{sourcepointer-4}{}}
\newlabel{lemma-immersion-permanence@cref}{{[originallemma][0][]07RK}{[1][0][]source}{}{}{}}
\makeatother
\crefname{originallemma}{original lemma}{original lemmas}
\Crefname{originallemma}{Original lemma}{Original lemmas}
\makeatletter
\newlabel{lemma-locally-finite-type-characterize}{{01T2}{source}{}{sourcepointer-5}{}}
\newlabel{lemma-locally-finite-type-characterize@cref}{{[originallemma][0][]01T2}{[1][0][]source}{}{}{}}
\makeatother
\crefname{originallemma}{original lemma}{original lemmas}
\Crefname{originallemma}{Original lemma}{Original lemmas}
\makeatletter
\newlabel{properties-lemma-invert-s-sections}{{01PW}{source}{}{sourcepointer-6}{}}
\newlabel{properties-lemma-invert-s-sections@cref}{{[originallemma][0][]01PW}{[1][0][]source}{}{}{}}
\makeatother
\crefname{originalproposition}{original proposition}{original propositions}
\Crefname{originalproposition}{Original proposition}{Original propositions}
\makeatletter
\newlabel{properties-proposition-characterize-ample}{{01Q3}{source}{}{sourcepointer-7}{}}
\newlabel{properties-proposition-characterize-ample@cref}{{[originalproposition][0][]01Q3}{[1][0][]source}{}{}{}}
\makeatother
\crefname{originallemma}{original lemma}{original lemmas}
\Crefname{originallemma}{Original lemma}{Original lemmas}
\makeatletter
\newlabel{schemes-lemma-composition-immersion}{{02V0}{source}{}{sourcepointer-8}{}}
\newlabel{schemes-lemma-composition-immersion@cref}{{[originallemma][0][]02V0}{[1][0][]source}{}{}{}}
\makeatother
\title{A locally finite type morphism with an ample invertible sheaf becomes a projective immersion after twisting by a large enough power}
\author{The Stacks Project Authors}
\date{stacks-project:morphisms.tex:local-snapshot-20261012; GNU FDL 1.2 or later}
\begin{document}
\maketitle
\noindent\textit{Editorial note.} All delivered mathematical text is the author's own, copied from the pinned Stacks Project chapter file pinned-original-source.tex; no statement or proof step was rewritten, reordered or removed.\par\medskip
\noindent\textit{Editorial note.} Following the multi-claim convention for this source, source\_url is the source-level URL https://stacks.math.columbia.edu/ and the source label lemma-quasi-projective-finite-type-over-S is carried in claim\_key and here rather than in a per-claim tag URL.\par\medskip
\noindent\textit{Editorial note.} The preamble is the chapter's shared Stacks Project preamble, retained as sources/stacks-morphisms/retained-preamble.tex: a line-for-line excerpt of the shared Stacks preamble staged with this snapshot. No macro was invented or redefined, and this card adds no preamble text of its own.\par\medskip
\noindent\textit{Editorial note.} Statement span 9412-9425 and proof span 9427-9515 come from the halfA pool card operations/pool-cards/stacks-morphisms-halfA.json and were re-read in the pinned chapter file before this card was written.\par\medskip
\noindent\textit{Editorial note.} The proof contains the author's aside that the first half proves the result 'for some d and n which is enough in virtually all applications' before the upgrade to all large powers; that sentence is kept verbatim.\par\medskip
\noindent\textit{Editorial note.} 8 cross-references occur inside the counted spans; they are listed in external\_references as navigation pointers only.\par\medskip
\noindent\textit{Allowed standard prerequisites (editorial).}\begin{itemize}
\item The characterization of ampleness by the basis of the topology given by the affine opens defined by homogeneous forms in the section ring. Cited at proof lines 9427-9432 and again at proof lines 9483-9485. Produces the finite cover by affine opens cut out by homogeneous sections and, later, the global generation of all large powers.
\item The description of sections of a power of an invertible sheaf over a distinguished affine open as quotients of homogeneous sections, and the characterization of locally finite type morphisms by finite type ring maps on affine charts. Cited at proof lines 9441-9449. Turns generators of the coordinate rings of the affine opens into homogeneous sections of a single large power.
\item The construction of a morphism to projective space from a generating set of sections of an invertible sheaf, and the identification of the pullback of the twisting sheaf. Cited at proof lines 9459-9462. Produces the morphism j with the correct pullback of the twisting sheaf.
\item Permanence of immersions under base change and products, the Segre embedding, and closedness of composites of immersions. Cited at proof lines 9508-9515. Shows that the morphism produced in the second half is again an immersion.
\end{itemize}\medskip
\ifdefined\hypertarget\hypertarget{sourcepointer-1}{}\fi
\noindent\textit{Source reference.} Original definition 01NF: focused-seeds/AGP/sources/stacks-snapshot-20261011/tags.txt, source lines 2156--2156. Same-chapter reference used in the retained author text at source line(s) 9465. It is supplied only as a navigation pointer so that the retained text compiles; the counted item delivers the author's own statement and the author's complete proof, and no part of the referenced item is copied in or substituted for the counted argument.\par\medskip
\ifdefined\hypertarget\hypertarget{sourcepointer-2}{}\fi
\noindent\textit{Source reference.} Original lemma 01NE: focused-seeds/AGP/sources/stacks-snapshot-20261011/tags.txt, source lines 2155--2155. Same-chapter reference used in the retained author text at source line(s) 9464. It is supplied only as a navigation pointer so that the retained text compiles; the counted item delivers the author's own statement and the author's complete proof, and no part of the referenced item is copied in or substituted for the counted argument.\par\medskip
\ifdefined\hypertarget\hypertarget{sourcepointer-3}{}\fi
\noindent\textit{Source reference.} Original lemma 01WD: focused-seeds/AGP/sources/stacks-snapshot-20261011/tags.txt, source lines 2478--2478. Same-chapter reference used in the retained author text at source line(s) 9508, 9689, 10652. It is supplied only as a navigation pointer so that the retained text compiles; the counted item delivers the author's own statement and the author's complete proof, and no part of the referenced item is copied in or substituted for the counted argument.\par\medskip
\ifdefined\hypertarget\hypertarget{sourcepointer-4}{}\fi
\noindent\textit{Source reference.} Original lemma 07RK: focused-seeds/AGP/sources/stacks-snapshot-20261011/tags.txt, source lines 9870--9870. Same-chapter reference used in the retained author text at source line(s) 361, 9510, 9691, 10764. It is supplied only as a navigation pointer so that the retained text compiles; the counted item delivers the author's own statement and the author's complete proof, and no part of the referenced item is copied in or substituted for the counted argument.\par\medskip
\ifdefined\hypertarget\hypertarget{sourcepointer-5}{}\fi
\noindent\textit{Source reference.} Original lemma 01T2: focused-seeds/AGP/sources/stacks-snapshot-20261011/tags.txt, source lines 2359--2359. Same-chapter reference used in the retained author text at source line(s) 2717, 2739, 2801, 2863, 3019, 3035, 3157, 3348, 3613, 3707, 3868, 5197, 5282, 6702, 9325, 9390, 9440, 11421, 13422, 13686. It is supplied only as a navigation pointer so that the retained text compiles; the counted item delivers the author's own statement and the author's complete proof, and no part of the referenced item is copied in or substituted for the counted argument.\par\medskip
\ifdefined\hypertarget\hypertarget{sourcepointer-6}{}\fi
\noindent\textit{Source reference.} Original lemma 01PW: focused-seeds/AGP/sources/stacks-snapshot-20261011/tags.txt, source lines 2245--2245. Same-chapter reference used in the retained author text at source line(s) 8858, 9446. It is supplied only as a navigation pointer so that the retained text compiles; the counted item delivers the author's own statement and the author's complete proof, and no part of the referenced item is copied in or substituted for the counted argument.\par\medskip
\ifdefined\hypertarget\hypertarget{sourcepointer-7}{}\fi
\noindent\textit{Source reference.} Original proposition 01Q3: focused-seeds/AGP/sources/stacks-snapshot-20261011/tags.txt, source lines 2252--2252. Same-chapter reference used in the retained author text at source line(s) 8676, 8876, 9431, 9484, 9655, 10874. It is supplied only as a navigation pointer so that the retained text compiles; the counted item delivers the author's own statement and the author's complete proof, and no part of the referenced item is copied in or substituted for the counted argument.\par\medskip
\ifdefined\hypertarget\hypertarget{sourcepointer-8}{}\fi
\noindent\textit{Source reference.} Original lemma 02V0: focused-seeds/AGP/sources/stacks-snapshot-20261011/tags.txt, source lines 3751--3751. Same-chapter reference used in the retained author text at source line(s) 180, 246, 9514, 9693. It is supplied only as a navigation pointer so that the retained text compiles; the counted item delivers the author's own statement and the author's complete proof, and no part of the referenced item is copied in or substituted for the counted argument.\par\medskip
\begin{lemma}
\label{lemma-quasi-projective-finite-type-over-S}
Let $f : X \to S$ be a morphism of schemes.
Let $\mathcal{L}$ be an invertible sheaf on $X$.
Assume that
\begin{enumerate}
\item the invertible sheaf $\mathcal{L}$ is ample on $X$, and
\item the morphism $X \to S$ is locally of finite type.
\end{enumerate}
Then there exists a $d_0 \geq 1$ such that for every $d \geq d_0$
there exist an $n \geq 0$ and an immersion
$i : X \to \mathbf{P}^n_S$ over $S$ such that
$\mathcal{L}^{\otimes d} \cong i^*\mathcal{O}_{\mathbf{P}^n_S}(1)$.
\end{lemma}

\begin{proof}
Let
$A = \Gamma_*(X, \mathcal{L}) =
\bigoplus_{d \geq 0} \Gamma(X, \mathcal{L}^{\otimes d})$.
By Properties, Proposition \ref{properties-proposition-characterize-ample}
the set of affine opens $X_a$ with $a \in A_{+}$ homogeneous forms
a basis for the topology of $X$. Hence we can find finitely
many such elements $a_0, \ldots, a_n \in A_{+}$ such that
\begin{enumerate}
\item we have $X = \bigcup_{i = 0, \ldots, n} X_{a_i}$,
\item each $X_{a_i}$ is affine, and
\item each $X_{a_i}$ maps into an affine open $V_i \subset S$.
\end{enumerate}
By Lemma \ref{lemma-locally-finite-type-characterize}
we see that the ring maps
$\mathcal{O}_S(V_i) \to \mathcal{O}_X(X_{a_i})$ are
of finite type. Hence we can find finitely many
elements $f_{ij} \in \mathcal{O}_X(X_{a_i})$, $j = 1, \ldots, n_i$
which generate $\mathcal{O}_X(X_{a_i})$ as an $\mathcal{O}_S(V_i)$-algebra.
By Properties, Lemma \ref{properties-lemma-invert-s-sections}
we may write each
$f_{ij}$ as $a_{ij}/a_i^{e_{ij}}$ for some
$a_{ij} \in A_{+}$ homogeneous. Let $N$ be a positive integer which
is a common multiple of all the degrees of the elements
$a_i$, $a_{ij}$. Consider the elements
$$
a_i^{N/\deg(a_i)}, \ a_{ij}a_i^{(N/\deg(a_i)) - e_{ij}} \in A_N.
$$
By construction these generate the invertible sheaf
$\mathcal{L}^{\otimes N}$ over $X$. Hence they give rise
to a morphism
$$
j : X \longrightarrow
\mathbf{P}_S^{m}
\quad
\text{with } m = n + \sum n_i
$$
over $S$, see Constructions, Lemma \ref{constructions-lemma-projective-space}
and Definition \ref{constructions-definition-projective-space}.
Moreover, $j^*\mathcal{O}_{\mathbf{P}_S}(1) = \mathcal{L}^{\otimes N}$.
We name the homogeneous coordinates $T_0, \ldots, T_n, T_{ij}$
instead of $T_0, \ldots, T_m$.
For $i = 0, \ldots, n$ we have $j^{-1}(D_{+}(T_i)) = X_{a_i}$.
Moreover, pulling back the element $T_{ij}/T_i$ via $j^\sharp$ we
get the element $f_{ij} \in \mathcal{O}_X(X_{a_i})$.
Hence the morphism $j$ restricted to $X_{a_i}$
gives a closed immersion of $X_{a_i}$ into the affine open
$D_{+}(T_i) \cap \mathbf{P}^m_{V_i}$ of $\mathbf{P}^N_S$.
Hence we conclude that the morphism $j$ is an immersion.
This implies the lemma holds for some $d$ and $n$ which is enough
in virtually all applications.

\medskip\noindent
This proves that for one $d_2 \geq 1$
(namely $d_2 = N$ above), some $m \geq 0$ there exists some
immersion $j : X \to \mathbf{P}^m_S$ given by global sections
$s'_0, \ldots, s'_m \in \Gamma(X, \mathcal{L}^{\otimes d_2})$.
By Properties, Proposition \ref{properties-proposition-characterize-ample}
we know there exists an integer
$d_1$ such that $\mathcal{L}^{\otimes d}$ is globally generated
for all $d \geq d_1$. Set $d_0 = d_1 + d_2$. We claim that
the lemma holds with this value of $d_0$. Namely, given
an integer $d \geq d_0$ we may choose $s''_1, \ldots, s''_t
\in \Gamma(X, \mathcal{L}^{\otimes d - d_2})$ which generate
$\mathcal{L}^{\otimes d - d_2}$ over $X$. Set $k = (m + 1)t$ and
denote $s_0, \ldots, s_k$ the collection of sections
$s'_\alpha s''_\beta$, $\alpha = 0, \ldots, m$,
$\beta = 1, \ldots, t$. These generate $\mathcal{L}^{\otimes d}$
over $X$ and therefore define a morphism
$$
i : X \longrightarrow \mathbf{P}^{k - 1}_S
$$
such that $i^*\mathcal{O}_{\mathbf{P}^n_S}(1) \cong \mathcal{L}^{\otimes d}$.
To see that $i$ is an immersion, observe that $i$ is the composition
$$
X \longrightarrow \mathbf{P}^m_S \times_S \mathbf{P}^{t - 1}_S
\longrightarrow \mathbf{P}^{k - 1}_S
$$
where the first morphism is $(j, j')$ with $j'$ given by
$s''_1, \ldots, s''_t$ and the
second morphism is the Segre embedding
(Constructions, Lemma \ref{constructions-lemma-segre-embedding}).
Since $j$ is an immersion, so is $(j, j')$
(apply Lemma \ref{lemma-immersion-permanence}
to $X \to \mathbf{P}^m_S \times_S \mathbf{P}^{t - 1}_S
\to \mathbf{P}^m_S$). Thus $i$ is a composition of
immersions and hence an immersion
(Schemes, Lemma \ref{schemes-lemma-composition-immersion}).
\end{proof}

\end{document}
